\section{Authors Note}
This area is new and information will be added to this document as route development continues. If you would like to contribute your knowledge to this project, please get in touch with me (I can be contacted via GitHub or my personal email: an.child@gmail.com). Every time this guidebook is updated the latest version is uploaded as a PDF to the project github page, follow the QR code below to download the latest version.
\setbox0=\hbox{\includegraphics[width=0.4\linewidth]{./images/maps/qr/Idleyld_qr.png}}% Store image in \box0
\needspace{\ht0}% Need at least the height of \box0
\begin{minipage}[c]{0.5\linewidth}{
\includegraphics[width=0.8\linewidth]{./images/maps/qr/Idleyld_qr.png}}
\end{minipage}
\begin{minipage}[c]{0.5\linewidth}{
\textcolor{blue}{\parbox{0.8\linewidth}{\href{https://github.com/AndrewChild/Idleyld_Guidebook}{\ul{Get the latest revision of this book}}}}}
\end{minipage}
\section{Navigation}
This book uses maps to show the layout of areas, additionally the approximate GPS locations of formations are provided in the index, GPS accuracy is poor in some places. A complete GPS map is available via Caltopo, this map is manually maintained and may not always be up to date.
\setbox0=\hbox{\includegraphics[width=0.4\linewidth]{./images/maps/qr/navigation.png}}% Store image in \box0
\needspace{\ht0}% Need at least the height of \box0
\begin{minipage}[c]{0.5\linewidth}{
\includegraphics[width=0.8\linewidth]{./images/maps/qr/navigation.png}}
\end{minipage}
\begin{minipage}[c]{0.5\linewidth}{
\textcolor{blue}{\parbox{0.8\linewidth}{\href{https://caltopo.com/m/B6P1P9M}{\ul{Caltopo map of formations}}}}}
\end{minipage}
\section{Amenities}
\subsection*{Toilets}
There are no toilets, bring a shovel and TP. Be sure to dig your cat hole well away from climbing and camping areas.
\subsection*{Camping}
The Parking area is on BLM land and camping is permitted. If you camp at the trail head be courteous to other visitors who will need to pass through your camp to access the climbing.
\subsection*{Water}
Bring your own.
\subsection*{Groceries}
The Idleyld store has a small assortment of produce and a deli along with the standard offerings that you might expect at a gas station. The nearest real grocery stores are about half an hour away in Roseburg. Visitors staying for multiple days might want to stock up on supplies before arriving.
\section{Local Ethics}
Most of the climbing is located on private timber land. The are currently no access issues and there is no relationship between the landowners and climbers. The established ethic for climbing on timber land in Oregon is that the owners prefer not to get involved, consequently climbers should do their best to keep a low profile and ensure the land owners don't need to get involved. There are a few specific activities which could threaten access for everyone:\\
\begin{itemize}
\item Building fires or causing fire hazards.\\
\item Parking on or blocking gated forest roads.\\
\item Failing to obey posted fire closures.\\
\end{itemize}
\subsection*{Open and Closed Projects}
This book notes several lines that have yet to see a first ascent. Generally boulder projects are understood to be "Open" which is to say that no one has dibs on them and they may be climbed by anyone. These projects are in fact included as a way to encourage and inspire climbers to check out and establish new lines.\\
\\
Rope climbs on the other hand may be either "Open" (anyone can get on it) or "Closed" (the developer has requested others to wait until they are finished establishing the route). Closed projects are customarily marked with a piece of red string or webbing on the first bolt.\\
There are many reasons why a route developer may choose to "Close" a project. It may be as simple as the developer hasn't finished cleaning and bolting the route. More commonly the developer is just requesting the privilege of the first ascent of the route and the naming rights that come along with it. Route development takes a good deal of time and money (each bolt on a sport climb costs upwards of \$6) thus it's considered reasonable for a route developer ask for a period of first dibs on the fruit of their labors. Failure to obey this request is considered route thievery and it's not a good way to make friends.\\
\section{Poison Oak}
Some sections of Idleyld are plagued with poison oak. Tread carefully and watch out for low growing shrubs with waxy leaves in clusters of three. The leaves turn red during the fall and fall off in the winter. Exposure to any part of the plant can cause irritation.\\
\begin{center}			  
\begin{overpic}[width=0.9\linewidth]{./images/photos/poison-oak.jpg}
\put (0,5) {\colorbox{\chapterColor}{\parbox{0.7\linewidth}{\textcolor{white}{Poison oak}}}}
\end{overpic}
\end{center}	
\section{Ticks}
Ticks are an increasing nuisance in Oregon due to our warming winter climate. It's common to encounter them while recreating at Idleyld in the spring and early summer. Be sure to check your self for tics after a day out.\\
\section{How to use this book}
\subsection*{Grades and Descriptions}
\colorbox{green!20}{\textbf{Boulder problems V0-V3}}\\
\colorbox{RoyalBlue!20}{\textbf{Boulder problems V4-V6}}\\
\colorbox{Goldenrod!50}{\textbf{Boulder problems V7-V9}}\\
\colorbox{red!20}{\textbf{Boulder problems V10+}}\\
\colorbox{green!20}{\textbf{Roped climbs 5.0-5.9}}\\
\colorbox{RoyalBlue!20}{\textbf{Roped climbs 5.10a-5.11d}}\\
\colorbox{Goldenrod!50}{\textbf{Roped climbs 5.12a-5.13d}}\\
\colorbox{red!20}{\textbf{Roped climbs 5.14a+}}\\
\colorbox{black!20}{\textbf{Projects and Unknown Grades}}\\
As much as possible the grades and descriptions of routes in this book have been based on the collective first hand experience of the collaborators of this book. Instances where first hand experience is limited or unavailable are graded with an asterisk.\\
\\
Boulder problems in this book are graded on the Hueco V scale and roped climbs are graded using the Yosemite decimal system. Although these grades are inherently subjective, care has been taken in considering the grading of each route. A color coding system is applied for ease of use as described at the beginning of this section.\\
\subsection*{Ratings for Quality and Seriousness}
In addition to a difficulty rating, route quality and seriousness ratings are provided on an out of three system as defined below.
\subsubsection{Quality}
\begin{tabular}{rcp{0.75\linewidth}}
&-&No quality rating given, this designation is typically only included for Projects and routes that the collaborators of this guide do not have first hand knowledge of.\\
\ding{72} \ding{72} \ding{72}&-&This route is an area classic, if you are unfamiliar with the area this is one you should check out on your first visit.\\
\ding{72} \ding{72}&-&This route is charming, but may be lacking one or more qualities of a true classic.\\
\ding{72}&-&This route may leave something to be desired but isn't objectively terrible.\\
\ding{73}&-&Zero stars, this route is bad.\\
\end{tabular}\\
\subsubsection{Seriousness}
\begin{tabular}{rcp{0.75\linewidth}}
&-&No seriousness rating given, this is generally a safe climb with appropriate padding/protection. There are no extraordinary hazards that you should be aware of.\\
\warn&-&A boulder with this rating may have insecure moves which are high off the ground or over a bad landing or both. A roped climb with this rating may have sections were falling presents risk of injury. A competent climber who is aware of these hazards will still be able to climb this at a minimally increased risk.\\
\warn \warn&-&There are sections of this climb where the risks are hard to minimize. Falls in certain areas may be unlikely for a climber of appropriate skill level but the consequences of such a fall could be real.\\
\warn \warn \warn&-&This route could cause serious injury or worse even when attempted by a person competent at climbing the assigned grade. This climb should be approached with caution.\\
\end{tabular}
